% questions.tex —— 1952 年·全国卷 · 题目内容片段（无答案）
% 由 tools/generate_projects.py 依来源总账自 1952/1952_数学真题_试题.md 生成，勿手改（改源 MD 后重跑）


\par\medskip

\gkpart{一、填空题}

\begin{question}
\gkti{1}分解因式：$x^4-y^4=$ \gkblank{}．

\end{question}
\begin{question}
\gkti{2}若 $\log_{10}2x=2\log_{10}x$，问 $x=$ \gkblank{}．

\end{question}
\begin{question}
\gkti{3}若方程式 $x^3+bx^2+cx+d=0$ 之三根为 $1$，$-1$，$\dfrac12$，则 $c=$ \gkblank{}．

\end{question}
\begin{question}
\gkti{4}若 $\sqrt{x^2+7}-4=0$，则 $x=$ \gkblank{}．

\end{question}
\begin{question}
\gkti{5}$\begin{vmatrix}1&2&3\\4&5&0\\3&2&1\end{vmatrix}=$ \gkblank{}．

\end{question}
\begin{question}
\gkti{6}两个圆的半径都是 $4$ 寸，并且一个圆通过另一圆的圆心，则这两个圆的公共弦之长是 \gkblank{} 寸．

\end{question}
\begin{question}
\gkti{7}三角形 $\triangle ABC$ 的面积是 $60$ 平方寸，$M$ 是 $AB$ 的中点，$N$ 是 $AC$ 的中点，则 $\triangle AMN$ 的面积是 \gkblank{} 平方寸．

\end{question}
\begin{question}
\gkti{8}正十边形的一内角是 \gkblank{} 度．

\end{question}
\begin{question}
\gkti{9}祖冲之的圆周率 $\pi=$ \gkblank{}．

\end{question}
\begin{question}
\gkti{10}球的面积等于大圆面积的 \gkblank{} 倍．

\end{question}
\begin{question}
\gkti{11}直圆锥之底之半径为 $3$ 尺，斜高为 $5$ 尺，则其体积为 \gkblank{} 立方尺．

\end{question}
\begin{question}
\gkti{12}正多面体有 \gkblank{} 种，其名称为 \gkblank{}．

\end{question}
\begin{question}
\gkti{13}已知 $\sin\theta=\dfrac13$，求 $\cos2\theta=$ \gkblank{}．

\end{question}
\begin{question}
\gkti{14}方程式 $\tan2x=1$ 的通解为 $x=$ \gkblank{}．

\end{question}
\begin{question}
\gkti{15}太阳仰角为 $30^\circ$ 时塔影长 $5$ 丈，求塔高 $=$ \gkblank{}．

\end{question}
\begin{question}
\gkti{16}三角形 $\triangle ABC$ 之 $b$ 边为 $3$ 寸，$c$ 边为 $4$ 寸，$A$ 角为 $30^\circ$，则 $\triangle ABC$ 的面积为 \gkblank{} 平方寸．

\end{question}
\begin{question}
\gkti{17}已知一直线经过点 $(2,-3)$，其斜率为 $-1$，则此直线之方程式为 \gkblank{}．

\end{question}
\begin{question}
\gkti{18}若原点在一圆上，而此圆的圆心为点 $(3,4)$，则此圆的方程式为 \gkblank{}．

\end{question}
\begin{question}
\gkti{19}原点至 $3x+4y+1=0$ 之距离 $=$ \gkblank{}．

\end{question}
\begin{question}
\gkti{20}抛物线 $y^2-8x+6y+17=0$ 之顶点之坐标为 \gkblank{}．

\par\medskip

\end{question}
\gkpart{二、解答题}

\begin{question}
\gkti{21}解方程式 $x^4+5x^3-7x^2-8x-12=0$．

\end{question}
\begin{question}
\gkti{22}$\triangle ABC$ 中，$\angle A$ 的外分角线与此三角形的外接圆相交于 $D$，求证：$BD=CD$．

\end{question}
\begin{question}
\gkti{23}设三角形的边长为 $a=4$，$b=5$，$c=6$，其对角依次为 $A$，$B$，$C$．

  (1) 求 $\cos C$；

  (2) 求 $\sin C$，$\sin B$，$\sin A$；

  (3) 问 $A$，$B$，$C$ 三个角各为锐角或钝角？

\end{question}
\begin{question}
\gkti{24}一椭圆通过 $(2,3)$ 及 $(-1,4)$ 两点，中心为原点，长短轴重合于坐标轴，试求其长短轴及焦点．

\end{question}